\scanpage{20}

\deleted{3) Si $\underline E$ est \deleted{de} de t.f. \inserted{et est quotient d'un Module loc libre (sans op.)} il existe un $\unsure{\underline O_X|\mathcal O_X}$-$G$-Module \deleted{\unsure{qui}} \uline{loc. libre de t.f.} $E$, et un épimorphisme de $\unsure{\underline O_X|\mathcal O_X}$-$G$-Modules $\underline F\longrightarrow E$.}

3) \deleted{Si $S$ est \unsure{et admet} un faisceau inv.} \inserted{qu. cpt \deleted{\unsure{quasi-sép.}} \inserted{qu. sép.} et si tt module quasi-coh.} de t.f. sur $S$ est $\simeq$ quotient d'un loc. libre de t.f., \deleted{ample,} alors $\underline E$ est $\simeq$ quotient d'un $\underline O_X$-$G$-Module loc. libre de type fini de la forme $f^*(H)(-n)$, \unsure{à}

\uline{Dém} Il suffit de prouver 1) et 3). Posons (où $f:X-S$)
\[
\underline{\Gamma}_{\deleted{*}}^\bullet(E)=\coprod_n f_*(E(n))
\]
Alors grâce au fait que $G$ est plat sur $S$ et $f$ qu. cpt et (qu) sép., on trouve que $\Gamma^\bullet(E)$ est un Module qu. coh sur $S$ sur lequel $G$ opère. Soit
\[
\underline{\Gamma}_{\deleted{*}}^\bullet(\unsure{\underline O_X|\mathcal O_X})=\deleted{\unsure{B^\bullet}}\mathcal A^\bullet
\]
c'est une Algèbre qu. coh sur laquelle $G$ opère. $E\unsure{\sim|\mapsto}\underline{\Gamma}^\bullet(E)$ est un foncteur en $E$, à valeurs dans les \deleted{\unsure{objets}} $\mathcal A^\bullet$-$G$-Modules qu. coh. Le foncteur est exact à gauche. On a un foncteur (\uline{exact})

\mspage{19}
